\noindent{\Large\bfseries Diagrama de clases del sistema Asistencia Motor 360}\\[0.1cm]
{\small Arquitectura Modelo--Vista--Controlador y estructuras de datos propias}\\[0.2cm]

\begin{tikzpicture}[x=1cm,y=1cm,scale=0.90,transform shape,
  clase/.style={draw=black!75,fill=white,rounded corners=1.5pt,align=left,
    font=\fontsize{5.15}{6.3}\selectfont,inner sep=2.5pt},
  iface/.style={clase,draw=blue!65,fill=blue!5},
  paquete/.style={draw=azulUD!55,dashed,fill=azulUD!4,rounded corners=4pt},
  etiquetaPaquete/.style={font=\fontsize{7}{8}\selectfont\bfseries,text=azulUD},
  rel/.style={-{Latex[length=1.5mm,width=1.2mm]},draw=black!70,semithick},
  dep/.style={-{Latex[length=1.5mm,width=1.2mm]},draw=black!60,dashed},
  comp/.style={-{Diamond[length=2mm,width=1.3mm]},draw=black!70,semithick},
  here/.style={font=\fontsize{5}{6}\selectfont,fill=white,inner sep=1pt,text=black!75}
]

% ============================ BANDAS ============================
\draw[paquete] (0.25,11.55) rectangle (27.25,18.55);
\node[etiquetaPaquete,anchor=south west] at (0.35,18.57) {VISTA (MVC)};
\draw[paquete] (0.25,9.15) rectangle (27.25,11.45);
\node[etiquetaPaquete,anchor=south west] at (0.35,11.47) {CONTROLADOR};
\draw[paquete] (0.25,2.50) rectangle (27.25,9.05);
\node[etiquetaPaquete,anchor=south west] at (0.35,9.07) {MODELO DEL DOMINIO};
\draw[paquete] (0.25,0.10) rectangle (27.25,2.40);
\node[etiquetaPaquete,anchor=north east] at (27.15,2.32) {ESTRUCTURAS DE DATOS PROPIAS};

% ============================ VISTA ============================
\node[iface,text width=3.0cm] (panelModulo) at (3.1,17.15) {%
\textbf{PanelModulo}\\[-1pt]
{\fontsize{4.8}{5.6}\selectfont \textit{«interface»}}\\
\rule{\linewidth}{0.25pt}\\
+ actualizar()};

\node[clase,text width=4.05cm] (ventana) at (3.6,14.55) {%
\textbf{VentanaPrincipal}\\
\rule{\linewidth}{0.25pt}\\
-- paneles: Lista<PanelModulo>\\
-- etiquetaEstado: JLabel\\
\rule{\linewidth}{0.25pt}\\
+ actualizarPaneles()\\
+ mostrarMensaje(String)};

\node[clase,text width=9.4cm] (paneles) at (14.1,16.35) {%
\textbf{Paneles de la interfaz} \textit{«vista»}\\
\rule{\linewidth}{0.25pt}\\
\begin{tabular}{@{}l@{\hspace{0.35cm}}l@{}}
PanelDashboard & PanelSolicitudes\\
PanelUnidades & PanelDespacho\\
PanelTecnicos & PanelInsumos\\
PanelClientes & PanelSupervisor\\
 & PanelReportes\\
\end{tabular}\\
\rule{\linewidth}{0.25pt}\\
+ actualizar() \hspace{0.2cm} \textit{implementan PanelModulo}};

\node[clase,text width=3.4cm] (modeloTabla) at (10.45,12.75) {%
\textbf{ModeloTabla}\\
{\fontsize{4.5}{5.4}\selectfont \textit{extends AbstractTableModel}}\\
\rule{\linewidth}{0.25pt}\\
-- filas: Lista<String[]>\\
+ getValueAt(int,int)\\
+ setFilas(Lista<String[]>)};

\node[clase,text width=3.5cm] (modeloCombo) at (14.25,12.75) {%
\textbf{ModeloCombo<T>}\\
{\fontsize{4.5}{5.4}\selectfont \textit{implements ComboBoxModel<T>}}\\
\rule{\linewidth}{0.25pt}\\
-- elementos: Lista<T>\\
+ setElementos(Lista<T>)\\
+ getSeleccionado(): T};

\node[clase,text width=3.65cm] (modeloLista) at (18.25,12.75) {%
\textbf{ModeloLista}\\
{\fontsize{4.5}{5.4}\selectfont \textit{extends AbstractListModel}}\\
\rule{\linewidth}{0.25pt}\\
-- elementos: Lista<String>\\
+ setElementos(Lista<String>)\\
+ getElementAt(int)};

% ======================= CONTROLADOR ===========================
\node[clase,text width=11.6cm] (controlador) at (13.6,10.30) {%
\textbf{ControladorPrincipal} \textit{«controlador»}\\
\rule{\linewidth}{0.25pt}\\
-- sistema: SistemaAsistencia \hspace{0.4cm} -- vista: VentanaPrincipal\\
\rule{\linewidth}{0.25pt}\\
+ registrarUnidad(...) \hspace{0.2cm} + registrarTecnico(...) \hspace{0.2cm} + registrarCliente(...)\\
+ registrarSolicitud(...) \hspace{0.2cm} + escalarSolicitud(String) \hspace{0.2cm} + despacharSiguiente(...)\\
+ despachoAutomatico() \hspace{0.2cm} + marcarAtendida(String) \hspace{0.2cm} + solucionarSolicitud(String)\\
+ deshacerUltimaTransaccion() \hspace{0.2cm} + exportarCSV(String) \hspace{0.2cm} + actualizarPaneles()};

% ============================ MODELO ===========================
\node[clase,text width=12.2cm] (sistema) at (14.5,7.65) {%
\textbf{SistemaAsistencia} \textit{«fachada del modelo»}\\
\rule{\linewidth}{0.25pt}\\
-- unidades, tecnicos, clientes, solicitudes: TablaHash\\
-- colaOrdinaria: Cola<Solicitud> \hspace{0.25cm} -- colaPrioritaria: ColaPrioridad<Solicitud>\\
-- serviciosActivos, serviciosSolucionados, historial: Lista\\
-- maletinesDisponibles, maletinesAuditoria, repuestos, transacciones: Pila\\
\rule{\linewidth}{0.25pt}\\
+ registrarUnidad(...) \hspace{0.15cm} + crearSolicitud(...) \hspace{0.15cm} + calcularPuntaje(...)\\
+ despachar(...) \hspace{0.15cm} + marcarAtendida(String) \hspace{0.15cm} + solucionarSolicitud(String)\\
+ deshacerUltimaTransaccion(): Transaccion \hspace{0.15cm} + exportarReporteCSV(String): int};

\node[clase,text width=4.3cm] (enums) at (24.2,7.95) {%
\textbf{Enumeraciones}\\
\rule{\linewidth}{0.25pt}\\
TipoUnidad, EstadoUnidad,\\
NivelDisponibilidad,\\
ZonaOperacion, Especialidad,\\
EstadoTecnico, TipoCliente,\\
TipoIncidente, EstadoSolicitud,\\
EstadoMaletin};

% ---- Entidades del dominio (fila inferior del modelo) ----
\node[clase,text width=2.45cm] (cliente) at (1.75,4.18) {%
\textbf{Cliente}\\
\rule{\linewidth}{0.25pt}\\
-- identificacion\\
-- nombre\\
-- tipo: TipoCliente\\
-- telefono\\
-- direccion};

\node[clase,text width=2.45cm] (unidad) at (4.45,4.18) {%
\textbf{UnidadMovil}\\
\rule{\linewidth}{0.25pt}\\
-- uuid: String\\
-- codigo: String\\
-- tipo: TipoUnidad\\
-- estado:\\
\hspace*{0.3cm}EstadoUnidad\\
-- zona, disponibilidad};

\node[clase,text width=2.45cm] (tecnico) at (7.15,4.18) {%
\textbf{Tecnico}\\
\rule{\linewidth}{0.25pt}\\
-- identificacion\\
-- nombre\\
-- especialidad\\
-- estado, sector\\
-- idServicioActivo};

\node[clase,text width=3.35cm] (solicitud) at (10.35,4.18) {%
\textbf{Solicitud}\\
\rule{\linewidth}{0.25pt}\\
-- id: String\\
-- cliente: Cliente\\
-- unidad: UnidadMovil\\
-- tecnico: Tecnico\\
-- tipoIncidente, zona\\
-- riesgo, afectacionVial\\
-- severidad\\
-- puntajePrioridad\\
-- prioritaria: boolean\\
-- estado: EstadoSolicitud\\
\rule{\linewidth}{0.25pt}\\
+ marcarAtendida()\\
+ marcarSolucionada()};

\node[clase,text width=2.45cm] (maletin) at (13.85,4.18) {%
\textbf{Maletin}\\
\rule{\linewidth}{0.25pt}\\
-- codigo\\
-- contenido\\
-- estado:\\
\hspace*{0.3cm}EstadoMaletin};

\node[clase,text width=2.45cm] (repuesto) at (16.55,4.18) {%
\textbf{Repuesto}\\
\rule{\linewidth}{0.25pt}\\
-- codigo\\
-- nombre};

\node[clase,text width=2.7cm] (transaccion) at (19.45,4.18) {%
\textbf{Transaccion}\\
\rule{\linewidth}{0.25pt}\\
-- consecutivo: int\\
-- descripcion: String\\
-- fecha: LocalDateTime\\
-- reversa: Runnable\\
\rule{\linewidth}{0.25pt}\\
+ deshacer()};

\node[clase,text width=2.45cm] (registro) at (22.3,4.18) {%
\textbf{RegistroAccion}\\
\rule{\linewidth}{0.25pt}\\
-- fecha\\
-- categoria\\
-- descripcion};

\node[clase,text width=2.75cm,font=\fontsize{4.7}{5.8}\selectfont] (exportador) at (25.1,4.18) {%
\textbf{ExportadorCSV}\\
\rule{\linewidth}{0.25pt}\\
+ exportar(Lista<Solicitud>,\\
\hspace*{0.45cm}String): int};

% ========================= ESTRUCTURAS =========================
\node[clase,text width=2.55cm] (lista) at (1.85,1.18) {%
\textbf{Lista<T>}\\
\rule{\linewidth}{0.25pt}\\
-- elementos: Object[]\\
-- tamano: int\\
\rule{\linewidth}{0.25pt}\\
+ agregar(T)\\
+ obtener(int): T\\
+ eliminar(int): T};

\node[clase,text width=2.45cm] (pila) at (5.0,1.18) {%
\textbf{Pila<T>}\\
\rule{\linewidth}{0.25pt}\\
-- elementos: Object[]\\
-- tamano: int\\
\rule{\linewidth}{0.25pt}\\
+ empujar(T)\\
+ desempujar(): T\\
+ cima(): T};

\node[clase,text width=2.45cm] (cola) at (8.1,1.18) {%
\textbf{Cola<T>}\\
\rule{\linewidth}{0.25pt}\\
-- frente, fin, tamano\\
\rule{\linewidth}{0.25pt}\\
+ encolar(T)\\
+ desencolar(): T\\
+ frente(): T};

\node[clase,text width=3.6cm] (colaPrioridad) at (11.25,1.18) {%
\textbf{ColaPrioridad<T>}\\
\rule{\linewidth}{0.25pt}\\
-- monticulo: Object[]\\
-- comparador: Comparador<T>\\
\rule{\linewidth}{0.25pt}\\
+ encolar(T)\\
+ desencolar(): T\\
+ mirar(): T};

\node[clase,text width=2.7cm] (tablaHash) at (14.75,1.18) {%
\textbf{TablaHash<K,V>}\\
\rule{\linewidth}{0.25pt}\\
-- cubetas:\\
\hspace*{0.3cm}NodoHash<K,V>[]\\
-- tamano: int\\
\rule{\linewidth}{0.25pt}\\
+ poner(K,V)\\
+ obtener(K): V\\
+ eliminar(K): V};

\node[iface,text width=2.4cm] (comparador) at (18.0,1.18) {%
\textbf{Comparador<T>}\\
{\fontsize{4.8}{5.6}\selectfont \textit{«interface»}}\\
\rule{\linewidth}{0.25pt}\\
+ comparar(T,T): int};

\node[clase,text width=2.6cm] (excepcion) at (21.2,1.18) {%
\textbf{EstructuraVacia}\\
\textbf{Excepcion}\\
\rule{\linewidth}{0.25pt}\\
\textit{extends RuntimeException}};

\node[clase,text width=4.2cm] (nodoHash) at (25.0,1.18) {%
\textbf{NodoHash<K,V>}\\
{\fontsize{4.8}{5.6}\selectfont \textit{clase interna de TablaHash}}\\
\rule{\linewidth}{0.25pt}\\
-- clave: K \hspace{0.2cm} -- valor: V\\
-- siguiente: NodoHash<K,V>};

% ========================= RELACIONES ==========================
\draw[dep] (paneles.west) -- node[here,above]{«implementan»} (panelModulo.east);
\draw[rel] (ventana.east) to[bend left=8] node[here,above]{1 \quad 1..*} (paneles.south west);
\draw[dep] (ventana.south) |- node[here,above,pos=0.25]{«usa»} (controlador.west);
\draw[dep] (paneles.east) -- (21.0,16.35) -- (21.0,10.30) -- (controlador.east);
\node[here,anchor=west] at (21.12,13.15) {«invoca»};
\draw[dep] (modeloCombo.north) -- node[here,right]{«usan los paneles»} (paneles.south);
\draw[dep] (controlador.south) -- node[here,right]{«delega operaciones»} (sistema.north);

\draw[comp] (sistema.south) -- (14.5,5.98);
\draw[black!70,semithick] (1.75,5.98) -- (25.1,5.98);
\foreach \n in {cliente,unidad,tecnico,solicitud,maletin,repuesto,transaccion,registro,exportador}
  \draw[black!70,semithick] (\n.north) -- (\n.north |- 0,5.98);

\draw[dep] (sistema.west) -- (0.35,7.65) -- (0.35,2.40);
\node[here,above] at (2.6,7.65) {«usa»};
\node[here,rotate=90] at (0.35,4.75) {estructuras propias};

\draw[dep] (colaPrioridad.north) .. controls (11.35,2.45) and (18.0,2.45) .. (comparador.north);
\node[here,above] at (15.0,2.52) {«usa»};

\end{tikzpicture}
